%%%%%%%%%%%%
%
% $Autor: Wings $
% $Datum: 2019-03-05 08:03:15Z $
% $Pfad: TemplateSensor $
% $Version: 4250 $
% !TeX spellcheck = en_GB/de_DE
% !TeX encoding = utf8
% !TeX root = filename 
% !TeX TXS-program:bibliography = txs:///biber
%
%%%%%%%%%%%%

% Structure
\chapter{Lichtschranke}
\section{Allgemein}




\section{Optische Näherungssensoren}

\subsection{Aufbau von Lichtschranken}


\subsection{Funktionsprinzip von Lichtschranken}





\section{Gabellichtschranke F249}
\section{Bibliothek}
\subsection{Beschriebung}

\subsection{Handbuch und Inbetriebnahme des Sensors}
Folgende Vorgehensweise kann zum Anschluss und Inbetriebnahme der Lichtschranke verwendet werden. 


\subsection{Anschlussbeispiel}



\subsection{Code}

%\subsection{Files}



\section{Kalibirerung und Justierung}


%\section{Simple Code}


\section{Anwendungsbeschreibung der Lichtschranke E18-D80NK}



\section{Tests}

\subsection{Einfacher Funktionstest}


\begin{center}
	\captionof{figure}{Dieser Code liest den digitalen Eingang 8 des Mikrocontrollers ein und erkennt anhand des Signals der Lichtschranke, ob ein Objekt den Lichtstrahl unterbricht. Der Zustand wird anschließend über die serielle Schnittstelle ausgegeben.} 
	\ArduinoExternal{}{../../Code/TestLichtschranke/TestLichtschranke.ino}
\end{center}


%\subsection{Test all Functions}
%
\section{Zusatzinformationen und weiterführende Literatur }


